\section{What kinematics studies}

\begin{definitionbox}
Kinematics is the part of mechanics that describes motion without explaining its causes.
\end{definitionbox}

At this stage we do not introduce forces, interactions, or causes of motion. Those questions belong to dynamics. In kinematics we ask a more basic question:
\[
\text{How can motion be described mathematically?}
\]

To answer this question, we use the mathematical language recalled in the previous chapter: coordinates, functions, parametric curves, vectors, derivatives, and integrals. The new idea is not a new mathematical object. The new idea is interpretation. We now use mathematical objects to describe the motion of a body in space and time.

The central object of kinematics is the position of a moving point as a function of time. Once position is known, we can ask how it changes. This leads to velocity. Once velocity is known, we can ask how velocity changes. This leads to acceleration:
\[
\text{position} \quad \longrightarrow \quad \text{velocity} \quad \longrightarrow \quad \text{acceleration}.
\]

Mathematically, this direction is built using derivatives. Later we will also move in the opposite direction:
\[
\text{acceleration} \quad \longrightarrow \quad \text{velocity} \quad \longrightarrow \quad \text{position}.
\]
This opposite direction is built using integrals and initial conditions.

\begin{remarkbox}
Kinematics describes motion. Dynamics explains why motion changes. Keeping these two questions separate is essential at the beginning.
\end{remarkbox}

This chapter should be read slowly. The goal is not only to remember formulas, but to understand how a mathematical description of motion is constructed.
